\subsection{Key Equations}
\RepeatEquation{eq:trace-map}

\begin{itemize}
    \item \textbf{Trace State}
        Trace Strains:
        \RepeatEquation{eq:def-gam-kap}
    \item Element
        \begin{description}
            \item[$\circ$ Exact] \RepeatEquation{eq:finite-alignment-field}
            \item[$\circ$ Linear]
            \RepeatEquation{eq:trace-lift-linear}
        \end{description}
    \item Section 
        \RepeatEquation{eq:saint-venant-stress}
        \RepeatEquation{eq:strain-tensor-decomp}
        \begin{description}
            \item[$\circ$ Wagner] \RepeatEquation{eq:wagner-gl-strain-pi}
            \item[$\circ$ Linear] \RepeatEquation{eq:linear-gl-strain-pi}
        \end{description}

    \item Saint~Venant, with \eqref{eq:trace-lift-linear}
     \begin{description}
        \item[Type 3]
          Trace Matrix:
          \RepeatEquation{eq:T0-inv}
          Fiber strain:
          \RepeatEquation{eq:eps-B-e}
          Resultants:
          \begin{description}
          \item[$\circ$ Petrov] 
          \RepeatEquation{eq:petrov-section}
          \item[$\circ$ Bubnov]
          \RepeatEquation{eq:bubnov-section}
          \end{description}
     \end{description}
\end{itemize}

\clearpage
\subsection{Summary}

\begin{itemize}
    \item A trace \(\mathcal{T}\) of a continuum displacement \(\hat{\bm u}\) is an operation:
    \RepeatEquation{eq:trace-map}
    whereby
    \RepeatEquation{eq:trace-lift-linear}
    \item Traces are physically distinguished by how they manifest in boundary fixities on the beam displacement \(\bm{u}\) and rotation \(\bm{\theta}\) fields.
    \item The geometric \citep{cowper1966shear} and energetic \citep{renton1991generalized} theories of shear  
     both represent exact continuum solutions, and are included in the Class I family of traces. 
    \item Any Class I trace (\cref{def:class-1}) is compatible with Type 3 beam finite elements via a Petrov-Galerkin section \eqref{eq:petrov-section}, which preserves the classical section resultants, but furnishes an asymmetric tangent stiffness.
    \item Class II traces (\cref{def:class-ii}) are the subset of Class I traces which may be implemented via a Bubnov-Galerkin section \eqref{eq:bubnov-section} to furnish a symmetric tangent stiffness.
    \item Class III traces (\cref{def:class-iii}) are the subset of Class II traces that produce the classical beam stress resultants \(\snm = (N,\ShearForce,T,\bm{M})\) under the symmetric Bubnov-Galerkin implementation \eqref{eq:bubnov-section}.
    \item All Class III traces are equivalent under Saint~Venant loads (\cref{def:sv-equivalence}), and their trace matrix is given by \eqref{eq:energetic-trace-blocks}.
\end{itemize}

\begin{itemize}
    \item There is a 36-parameter space of particular solutions to Saint Venant's problem that can be consistently modeled by standard beam elements with only enhancements to the resultant model.
    \item An 18 parameter subset of these solutions will admit a physically meaningful kinematic interpretation in the configuration variables \(\bm{u},\bm{\theta}\).
    \item A  4 parameter subset is exactly variationally compatible with the linear Type 3 equilibrium operator.
    \item There is a unique model.
\end{itemize}

Novelty:
\begin{itemize}
    \item Classical trace equations are obtained in a form suitable for 3D inelastic simulations with Type 3 elements, general sections, and non-zero Poisson ratio. 
    \item Identify exact scaling relationship between shear warping functions \(\bm{\psi}\) and beam strains \(\bm\gamma\).
    \item Nonlinear section geometry (Wagner term) incorporated inside section.
    \item Classical shear correction coefficients are associated to continuum exact 3D displacements.
\end{itemize}


TODO:
\begin{itemize}
    \item Implement Class I trace blocks.
    \item Implement trace blocks inside section.
    \item How stable are Class I traces?
    \item How useful are non-Class III traces?
    \item Consolidate trace matrix blocks.
    \item Investigate trace block relationships to section geometry.
\end{itemize}

Future:
\begin{itemize}
    \item Investigate traces incorporating material shear stress into Euler type beam elements.
    \item Investigate Class III traces outside of Saint~Venant loading.
\end{itemize}


\begin{Aside}

When one includes shear warping, the plane Timoshenko mode that classically scales with $\bm\gamma$ becomes parasitic. 

ie, in the saint venant solution, shear strains arise like 
$$
\bm\varepsilon_{\gamma}(x,\bm{r}) = (\nabla\bm{\WarpShearIesan}(\bm{r}))^{\top} X \bm\gamma(x)
$$ 
where $X$ is some constant ``shear correction", which is not yet exactly the classical shear correction coefficient, but eventually will give rise to it. 
This is incompatible with a naive enrichment of Timoshenko kinematics: 
$$\hat{\bm u}(x,\bm{r}) = \bm{u}(x) + \bm{\theta}(x)\times \bm{r} + warping$$
which would produce 
$$
\bm\varepsilon = \bm\gamma(x) + \bm\kappa(x)\times \bm{r} + warping-strain(x,r)
$$
Note the parasitic ``plane" $\bm{r}$-constant strain $\bm\gamma$. 
We formalize this situation by introducing the kinematic ansatz:
$$
\hat{\bm u}(x,r) = \bm u(x) + \bm\theta(x)\times \bm{r} 
+ \bm\Phi(\bm{r})\bm\alpha(x) + \bm\Psi(\bm{r})\bm\eta(x)
$$
where $\bm\Phi$ is a general matrix basis of section warping functions, $\bm\alpha(x)$ their amplitudes, and $\bm\Psi(\bm{r})$ reintroduces the section-rigid patterns that are also represented by $\bm{u}(x),\bm\theta(x)$. 
The auxiliary fields $\bm\alpha(x)$ and $\bm\eta(x)$ will eventually be eliminated to furnish a classical 6-DOF theory in only $\bm{u}(x),\bm\theta(x)$. 
\end{Aside}



Under the energy-conjugate trace, a clamped boundary condition 
$\bm{u}(0) = \bm{\theta}(0) = \mathbf{0}$ does not fix a direct geometric average over the section, but rather constrains the energy-conjugate kinematic measures in a way that may be difficult to realize physically. 
An illuminating geometric derivation of this model was put forth by \cite{serpieri2014equivalence}, which reveals it's interpretation as a long-wavelength approximation to the continuum problem of a tip-loaded cantilever that is \emph{fully} restrained at it's root \cite{ladeveze1998new}. 
This problem does not fall under the relaxed Saint~Venant class, which precludes Dirichlet conditions of this nature. 





















% 1. are these conclusions theoretically sound
% 2. Am i missing any salient points of practical interest?
% 3. How does one understand the implications of the boundary condition phenomenon on beam element continuity? does trace selection have an impact on compatibility in finite element analysis? eg, if i model a saint-venant transversely loaded cantilever with two beam finite finite elements under a given trace, what does the displacement field look like at the element interface?